\begin{Remark}\label{rem:MoritaRKK}
There is a natural notion of Morita equivalence of $C^*$-algebra bundles over~$X$, similar to the notion of
equivariant Morita equivalence from Theorem~\ref{equivMorita}, which
uses the $C_0(X)$-bimodule structures: a~\emph{$C_0(X)$-linear Morita equivalence} between two
$C_0(X)$-algebras is a Morita equivalence which is compatible with the
$C^*$-bundle structures over~$X$. More generally, there is a category
$\mathcal{R}\CCK\hspace{-1mm}\CCK_X$ of $C^*$-algebra bundles over $X$ whose morphisms are
elements of Kasparov's groups $\mathcal{R}{\rm KK}(X;\alg,\balg)$, see, e.g., \cite{Echterhoff2009}: the cycles are the usual cycles
$(\hil,\phi,T)$ for Kasparov's bivariant K-theory ${\rm
 KK}(\alg,\balg)$ (cf.\ Section~\ref{sec:TdualityKK}) with the
additional requirement that $\phi\colon \alg\to\End_\balg(\hil)$ is
$C_0(X)$-linear. There is an obvious faithful functor
$\mathcal{R}\CCK\hspace{-1mm}\CCK_X \to \CCK\hspace{-1mm}\CCK$ which forgets the
$C_0(X)$-algebra structures. Isomorphic $C^*$-bundles in the category
$\mathcal{R}\CCK\hspace{-1mm}\CCK_X$ are precisely the $\mathcal{R}{\rm KK}$-equivalent
$C^*$-bundles. If $\alg$ and $\balg$ are isomorphic in
$\mathcal{R}\CCK\hspace{-1mm}\CCK_X$, i.e., there exists an invertible class
$\alpha\in \mathcal{R}{\rm KK}(X;\alg,\balg)$, then they are also
isomorphic as
$C^*$-algebras in the category $\CCK\hspace{-1mm}\CCK$.
\end{Remark}

\subsection[$\real^n$-actions on Mostow bundles]{$\boldsymbol{\real^n}$-actions on Mostow bundles}\label{sec:RMostow}

We can now apply the results of Section~\ref{sec:TopTtorus} to the
class of twisted tori given in
Section~\ref{sec:abeliansolv}.
The Mostow fibration of any almost abelian solvmanifold identifies
$\torus_{\Lambda_\sfG}$ as a torus bundle over a~circle, hence the
algebra of functions $C(\torus_{\Lambda_\sfG})$ is a~$C(\torus)$-algebra. In other words, $C(\torus_{\Lambda_\sfG})$ is an
object of the category $\mathcal{R}\CCK\hspace{-1mm}\CCK_\torus$, and we are
interested in the T-duality isomorphisms of~$C(\torus_{\Lambda_\sfG})$
in this category. In particular, given a fibrewise right action of the abelian
Lie group~$\R^n$ on~$\torus_{\Lambda_\sfG}$, it follows from
Theorem~\ref{thm:crossedfibres} that the $C^*$-algebraic T-dual
$C(\torus_{\Lambda_\sfG})\rtimes_\rt\R^n$ is also a $C(\torus)$-algebra, and by~\cite[Theorem~3.5]{Echterhoff2009} the $C^*$-bundles
$C(\torus_{\Lambda_\sfG})$ and
$C(\torus_{\Lambda_\sfG})\rtimes_\rt\R^n$ are isomorphic as $C^*$-algebras in the
category $\mathcal{R}\CCK\hspace{-1mm}\CCK_\torus$.

In order to have sensible definitions of
T-duality, we need to identify the
homologically non-trivial one-cycles of the twisted torus
$\torus_{\Lambda_\sfG}$, which are determined
in~\cite[Proposition~4.7]{Bock2009}. Write
\[
\ttM = (m_{ij})
\]
for the
integer matrix elements $m_{ij}\in\Z$ of the monodromy matrix. Since
$\sfG=\real^{d-1}\rtimes_\varphi\real$ is simply-connected, the
fundamental group of the twisted torus is
$\pi_1(\torus_{\Lambda_\sfG})\simeq\Lambda_\sfG$ whose abelianisation
$\Lambda_\sfG/[\Lambda_\sfG,\Lambda_\sfG]$ gives
the first homology group via the presentation
\begin{gather}\label{eq:H1torus}
{\rm H}_1(\torus_{\Lambda_\sfG},\zed) = \zed \oplus \bigg\langle
\hat e_1,\dots,\hat e_{d-1} \, \Big| \, \sum_{j=1}^{d-1} m_{ji} \hat
e_j = \hat e_i
\ \mbox{for} \
i=1,\dots,d-1 \bigg\rangle,
\end{gather}
where the first factor of $\zed$ corresponds to the base circle of the
torus fibration and the generators $\hat e_1,\dots,\hat e_{d-1}$
correspond to the torus fibres; they are given by
\begin{gather}\label{eq:hatei}
\hat e_i:=\sum_{k=1}^{d-1} \Sigma_{ki} \vec e_k,
\end{gather}
where
$\vec e_1,\ldots, \vec e_{d-1}$ is the standard basis of $\zed^{d-1}$ giving the
group law~\eqref{eq:grouplaw}.

We can then apply the structure theorem for finitely-generated
$\zed$-modules by appealing to some classical matrix algebra.
From the presentation \eqref{eq:H1torus}, ${\rm
 H}_1(\torus_{\Lambda_\sfG},\zed) = \zed \oplus {\rm{coker}}(\varphi_{x_0}-{\rm{id}}_{\Z^{d-1}})$
where $(\varphi_{x_0}-{\rm{id}}_{\Z^{d-1}})\colon \zed^{d-1}\to\zed^{d-1}$ in the
basis $\hat e_1,\ldots,\hat e_{d-1}$ is given by the integer relation matrix
\begin{gather}\label{eq:relationmatrix}
{\tt A}:=\ttM-\unit_{d-1}~.
\end{gather}
Let $r$ be the rank of $\tt A$. This matrix can be brought into its Smith normal form~$\tt D$
by finding invertible integer matrices ${\tt L}, {\tt
 R}\in\sfGL(d-1,\zed)$ such that
\[
{\tt D} = {\tt L} {\tt A} {\tt R}
\]
is diagonal with entries $m_i\in\zed$ for
$i=1,\dots,d-1$. The
integers~$m_i$ are the \emph{elementary divisors} of
${\tt A}$. They have the properties that $m_i$ divides
$m_{i+1}$, for $0<i<d-1$, and in particular $m_i=0$ for $i>r$; they can
be computed explicitly (up to sign) as
\[
m_i = \frac{d_i({\tt A})}{d_{i-1}({\tt A})},
\]
where the \emph{$i$-th determinant divisor} $d_i({\tt A})$ is the greatest common
divisor of all $i{\times} i$ minors of the relation matrix ${\tt A}$,
with $d_0({\tt A}):=1$. The matrices ${\tt L}, {\tt
 R}\in\sfGL(d-1,\zed)$ are found by reducing the matrix ${\tt A}$ to
its Smith normal form ${\tt D}$ through a sequence of elementary row
and column operations over~$\zed$, see, e.g.,~\cite{Havas1997}.

Given the Smith normal form, we set
\begin{gather}\label{eq:tildeei}
\tilde e_i:=\sum_{j=1}^{d-1}\,\big({\tt L}^{-1}\big)_{ji}\,\hat e_j
\end{gather}
and observe that
the image of $\tt A$, which is generated over $\zed$ by the vectors
\[
\sum_{j=1}^{d-1} {\tt A}_{ji} \hat e_j=\sum_{k,l=1}^{d-1} \big({\tt
 R}^{-1}\big)_{ki} {\tt
 D}_{lk} \tilde e_k ,
\]
is equivalently generated by the
vectors $m_k \tilde e_k$ with $k=1,\ldots, r$.
Hence
\begin{align*}
{\rm{coker}}(\varphi_{x_0}-{\rm{id}}_{\Z^{d-1}})&=\big\langle \hat e_1, \ldots , \hat
e_{d-1}\big\rangle\Big/\bigg\langle \sum\limits_{j=1}^{d-1} {\tt A}_{j1} \hat e_j,\ldots,
 \sum\limits_{j=1}^{d-1} {\tt A}_{j\:\!d-1} \hat e_j\bigg\rangle\\
&=\big\langle \tilde e_1,
\ldots , \tilde e_{d-1}\big\rangle\big/\big\langle m_1 \tilde e_1,\ldots, m_r \tilde
e_r\big\rangle
\end{align*}
and
\begin{gather}\label{eq:H1gendecomp}
{\rm H}_1(\torus_{\Lambda_\sfG},\zed) \simeq \zed \oplus
\zed^{d-1-r} \oplus \bigoplus_{i=1}^r \zed_{m_i} .
\end{gather}

We are exclusively
interested in the natural $\R^n$-actions on $\torus_{\Lambda_\sfG}$ which descend
from actions of abelian subgroups $\R^n\subset\sfG$, acting on $\sfG$
by right multiplication. They can be
organised into three classes associated with the different types of
summands in the $\Z$-module presentation of the homology group~\eqref{eq:H1gendecomp}, and we only retain those which are fiberwise
actions on the Mostow bundle.
The first summand~$\Z$ in~\eqref{eq:H1gendecomp} corresponds to the subgroup
\[
\R_x = \big\{(0,\xi)\in\sfG\big\}
\]
acting on $\sfG$ by right multiplication:
\[
(z,x) (0,\xi) = (z,x+\xi)
\]
for all $(z,x)\in\sfG$ and $\xi\in\R$. Clearly this does not descend
to a fiberwise
action on the twisted torus $\torus_{\Lambda_\sfG}$, and the crossed
product $C(\torus_{\Lambda_\sfG})\rtimes_\rt\R_x$ is no longer a
$C(\torus)$-algebra. Thus an $\R_x$-action takes us out of the
category $\mathcal{R}\CCK\hspace{-1mm}\CCK_\torus$, and we will henceforth
discard $\R^n$-actions where $\R^n$ contains the subgroup~$\R_x$.

For the remaining types of summands in \eqref{eq:H1gendecomp}, we can
give explicit descriptions of the $C^*$-algebraic T-duals of an almost
abelian solvmanifold. We consider both classes in turn. As the only solvmanifolds in one and two dimensions are tori, which are
already treated by our analysis from
Section~\ref{sec:toruscrossedprod}, we assume $d\geq3$ for the
remainder of this paper.

